\begin{abstract}
Tourism-dependent economies face substantial risks when international visitor flows are disrupted, yet evidence on the short-run labor market effects of sudden tourism supply shocks remains limited. This paper exploits the unexpected bankruptcy of WOW Air, a low-cost airline, in 2019 as a natural experiment that sharply curtailed international tourism to Iceland. Using a Bayesian Structural Time Series model, I estimate labor market outcomes relative to a counterfactual in which the airline did not exit. I find that WOW Air’s collapse reduced international passenger arrivals by more than 20\% and tourism employment by 8\% over the subsequent nine months, relative to the counterfactual, including a 4\% year-over-year decline in tourism jobs. Excluding air transport, tourism employment declined by 6.6\% relative to the counterfactual. At the national level, unemployment was nearly 50\% higher than would otherwise have been expected. These results highlight the vulnerability of tourism-dependent economies to sudden transport disruptions.
\\ 
\\
\normalfont
\noindent\textbf{Keywords:}  tourism employment; air transport; tourism supply shock; economic vulnerability; counterfactual analysis;  natural experiment
\end{abstract}
